\section{A 21-variable generator of degree 25}
\label{sec:var21}

At the other end of the frontier the game is the number of variables, and the
tool is elimination. We report two points, and we present the weaker one first
because its guarantee is stronger.

\subsection{Three eliminations with no side conditions: \texorpdfstring{$(23,25)$}{(23,25)}}

Three of the twenty-five unknowns of \cite{jsww1976} are determined linearly by
an equation whose right-hand side has all coefficients nonnegative:
\begin{align}
  z &= (gk+2g+k+1)(h+j)+h, \\
  q &= wz+h+j, \\
  y &= n+l+v .
\end{align}
By Remark~\ref{rem:crit} each may be substituted away with no argument beyond
the shape of the expression, and by Proposition~\ref{prop:elim} the resulting
system in $22$ unknowns is equisatisfiable with the original. The largest degree
is unchanged at $12$, so

\begin{theorem}
\label{thm:23}
There is a generator of the primes with $\nu=23$ and $\delta=25$.
\end{theorem}

This is three variables below the $(26,25)$ of \cite{jsww1976}, at the same
degree, and it uses no number theory whatsoever: it is three substitutions. That
it had not been written down is, we think, evidence that this part of the
frontier is less well explored than its age suggests.

\subsection{Two more, at the cost of a Pell bound: \texorpdfstring{$(21,25)$}{(21,25)}}

Two further unknowns, $m$ and $x$, are linearly determined, but with mixed
signs: their right-hand sides encode inequalities and
Remark~\ref{rem:crit} does not apply. Both become eliminable once one knows
\begin{equation}
  \label{eq:pell-bound}
  a \;\ge\; e+1 ,
\end{equation}
where $e=2n+p+q+z$. Rather than assume \eqref{eq:pell-bound}, we make it
structural by the substitutions $n=N+2$ and $a=e+1+A$ with $N,A\ge0$, which is a
bijective reparametrization of the region \eqref{eq:pell-bound} carves out. In
the reparametrized system all five right-hand sides have nonnegative
coefficients.

\begin{theorem}
\label{thm:21}
Assume \eqref{eq:pell-bound}. Then there is a generator of the primes with
$\nu=21$ and $\delta=25$.
\end{theorem}

\subsection{The bound \texorpdfstring{\eqref{eq:pell-bound}}{a >= e+1}, proved rather than cited}

Equation~(5) of the system reads $e^3(e+2)(a+1)^2+1=o^2$, and the factorization
\[
  e^3(e+2)\;=\;e^2\bigl((e+1)^2-1\bigr)
\]
turns it into the Pell equation $o^2-(A^2-1)Z^2=1$ with $A=e+1$ and $Z=e(a+1)$.
Three standard facts about that equation then give \eqref{eq:pell-bound}:

\begin{enumerate}[label=(P\arabic*),nosep]
  \item \emph{completeness}: every solution with $o,Z\ge0$ is $(X_j,Y_j)$ for
        some $j$, where $X_0=1$, $Y_0=0$ and
        $(X_{j+1},Y_{j+1})=(AX_j+(A^2-1)Y_j,\;X_j+AY_j)$;
  \item \emph{congruence}: $Y_j\equiv j \pmod{A-1}$;
  \item \emph{growth}: $Y$ is strictly increasing.
\end{enumerate}

From (P1), $Z=e(a+1)=Y_j$ for some $j$; from (P2), $e\mid Z$ forces $e\mid j$,
and $j\ge1$ gives $j\ge e$; from (P3) and $e\ge4$ we get
$Y_j\ge Y_3=4A^2-1$, whence $e(a+1)\ge e(e+2)$ and, cancelling $e>0$,
inequality~\eqref{eq:pell-bound}. The condition $e\ge4$ follows from $n\ge2$,
which is itself forced by equation~(4) by a squeeze between consecutive squares.

These three facts are available in the standard libraries of several proof
assistants; we prove all three from scratch (Section~\ref{sec:lean}) so that the
figure does not rest on a citation.

\subsection{Comparison with the literature, stated carefully}

\begin{center}
\begin{tabular}{@{}llll@{}}
\toprule
\textbf{Pair} & \textbf{Source} & \textbf{Status} & \textbf{Relation to ours} \\
\midrule
$(26,25)$ & \cite{jsww1976}, system (1) & exhibited & we improve by $3$, resp.\ $5$ \\
$(24,37)$ & Matiyasevich, via \cite{pak2022jsww} & announced & incomparable \\
$(21,21)$ & Matiyasevich (addendum) & announced & \textbf{dominates our $(21,25)$} \\
$(19,29)$ & \cite{jsww1976}, p.~450 & announced & incomparable \\
$(12,13697)$ & \cite{jsww1976}, Thm.~2 & announced & fewer variables, far higher degree \\
$(10,>6000)$ & \cite{matiyasevich1981} & exhibited & fewer variables and lower degree \\
\bottomrule
\end{tabular}
\end{center}

We are explicit about what this table says. Against the only \emph{exhibited}
generator at degree~$25$, our points improve the variable count by three and by
five. Against \emph{announced} figures they do not: the $(21,21)$ attributed to
Matiyasevich dominates our $(21,25)$, and the $(10,>6000)$
of~\cite{matiyasevich1981}---which is exhibited, and has been formalized in
Mizar \cite{pak2022ten}---dominates every point we obtain below $13$ variables.
Our contribution here is not a record in the number of variables; it is that
two points at the exhibited end of that axis now exist as checkable objects.
